\chapter{Absence d'arbitrage et théorèmes fondamentaux}\label{ch:dv:no-arbitrage-and-the-fundamental-theorems}

Quatre options sur l'indice S\&P~500 figurent dans un même ordre~: un call acheté et un call vendu,
un put vendu et un put acheté, aux strikes 5\,000 et 6\,000, toutes venant à échéance dans
un an. Quelle que soit l'évolution de l'indice, les quatre ensemble versent exactement 1\,000 points d'indice
à l'échéance, soit USD~100\,000 par ensemble de contrats. Le package se traite à
958.90 points. Personne sur le desk ne le considère comme une position d'options~: c'est une
obligation zéro-coupon à un an, garantie par la chambre de compensation, et 958.90 correspond à un
taux d'intérêt de 4.20~\%. Certains fonds prêtent des liquidités de cette façon, d'autres en empruntent~; la
bourse publie la taille du marché. Derrière cette transaction se trouve l'idée sur laquelle
repose chaque chapitre de ce livre~: deux portefeuilles au même paiement
doivent avoir le même prix, faute de quoi quelqu'un empoche la différence sans rien risquer. Ce
chapitre précise cette idée, la traduit en \omterm{def:dv:no-arbitrage-and-the-fundamental-theorems:ad}{prix d'états} du monde,
et démontre les deux théorèmes qui disent quand un prix est imposé et quand il n'est
que borné.

\section{Arbitrage et loi du prix unique}

Fixons une date $0$, une date $T$ et un ensemble fini d'états $\omega_1,\dots,\omega_S$,
dont l'un se réalisera en $T$. Il existe $N$ actifs négociés. L'actif $j$ coûte
$p_j$ aujourd'hui et verse $D_{sj}$ dans l'état $s$~: le marché est décrit par le vecteur de prix
$p\in\R^N$ et par la matrice de paiements de taille $S\times N$ notée $D$. Un portefeuille
$\theta\in\R^N$ (coefficients positifs pour les achats, négatifs pour les ventes) coûte $p\cdot\theta$
et verse le vecteur $D\theta$.

\begin{definition}[Actif contingent et réplication]\label{def:dv:no-arbitrage-and-the-fundamental-theorems:claim}
Un \emph{actif contingent}\index{actif contingent} est un paiement en $T$ qui
dépend de l'état~: un vecteur $g\in\R^S$ dans le modèle à une période, une variable aléatoire
$g(S_T)$ ou $g(S_t, t\le T)$ en général. Un \emph{portefeuille de
réplication}\index{portefeuille de réplication} pour $g$ est un portefeuille d'actifs négociés
dont le paiement est égal à $g$ dans chaque état~: $D\theta = g$.
\end{definition}

\begin{definition}[Arbitrage]\label{def:dv:no-arbitrage-and-the-fundamental-theorems:arbitrage}
Un \emph{arbitrage}\index{arbitrage} est un portefeuille qui ne coûte rien ou
coûte moins que rien, ne verse jamais un montant négatif, et n'est pas identiquement nul~: $p\cdot
\theta\le0$, $D\theta\ge0$, et au moins l'une des inégalités est stricte ($p\cdot
\theta<0$, ou $D\theta\ne0$). Le marché est \emph{sans arbitrage} lorsqu'aucun tel
portefeuille n'existe.
\end{definition}

\begin{definition}[Loi du prix unique]\label{def:dv:no-arbitrage-and-the-fundamental-theorems:loop}
La \emph{loi du prix unique}\index{loi du prix unique} est vérifiée lorsque deux portefeuilles
au même paiement ont le même coût~: $D\theta=D\theta'$ implique $p\cdot
\theta=p\cdot\theta'$. Sous cette hypothèse, tout actif réplicable a un prix unique,
le coût de n'importe lequel de ses \omterm{def:dv:no-arbitrage-and-the-fundamental-theorems:claim}{portefeuilles de réplication}.
\end{definition}

L'absence d'\omterm{def:dv:no-arbitrage-and-the-fundamental-theorems:arbitrage}{arbitrage} implique la \omterm{def:dv:no-arbitrage-and-the-fundamental-theorems:loop}{loi du prix unique} (si deux portefeuilles de même paiement
avaient des coûts différents, acheter le moins cher et vendre le plus cher serait un \omterm{def:dv:no-arbitrage-and-the-fundamental-theorems:arbitrage}{arbitrage}),
mais la réciproque est fausse~: un marché où un actif versant $(1,0)$
coûte zéro vérifie la \omterm{def:dv:no-arbitrage-and-the-fundamental-theorems:loop}{loi du prix unique} et offre pourtant un gain sans risque
(\cref{exo:dv:no-arbitrage-and-the-fundamental-theorems:4}). L'\omterm{def:dv:no-arbitrage-and-the-fundamental-theorems:arbitrage}{arbitrage}
est la notion la plus forte et la plus utile~: un desk qui en repère un l'exécute
en volume jusqu'à ce que les prix bougent.

\begin{definition}[Box spread]\label{def:dv:no-arbitrage-and-the-fundamental-theorems:box}
Un \emph{box spread}\index{box spread} de strikes $K_1<K_2$ et d'une même échéance est
un call acheté et un put vendu en $K_1$, auxquels s'ajoutent un call vendu et un put acheté en $K_2$.
À l'échéance, il verse $K_2-K_1$ dans chaque état, donc son prix vaut
$P(0,T)(K_2-K_1)$ et il se comporte comme une obligation zéro-coupon.
\end{definition}

Les deux premières jambes forment un terme long synthétique en $K_1$, les deux dernières un terme court
synthétique en $K_2$~; la différence de deux termes sur le même
actif est un montant certain (\cref{fig:dv:no-arbitrage-and-the-fundamental-theorems:box}).
L'argument repose sur la parité call-put (livre~1, chapitre~25) appliquée
deux fois, et, comme la parité, il exige un exercice européen~: avec des options américaines, les
jambes vendues peuvent être assignées par anticipation et le box cesse d'être une obligation.

\begin{omfigure}
\begin{tikzpicture}
\begin{axis}[width=11cm,height=5cm,
  xlabel={indice à l'échéance},ylabel={paiement (points d'indice)},
  tick label style={font=\small},label style={font=\small},
  xmin=4000,xmax=7000,ymin=-2100,ymax=2100,grid=major,grid style={gray!25},
  xtick={4000,5000,6000,7000},xticklabel style={/pgf/number format/.cd,1000 sep={\,}},
  yticklabel style={/pgf/number format/.cd,1000 sep={\,}},
  legend columns=3,legend style={font=\footnotesize,at={(0.5,-0.3)},anchor=north,draw=none,fill=none,/tikz/every even column/.append style={column sep=8pt}},
  table/col sep=comma]
\addplot[omDef,thick] table[x=s,y=longcall]{figdata/derivatives/01-no-arbitrage-and-the-fundamental-theorems/box_payoff.csv};
\addplot[omDef,thick,dashed] table[x=s,y=shortcall]{figdata/derivatives/01-no-arbitrage-and-the-fundamental-theorems/box_payoff.csv};
\addplot[omProp,thick] table[x=s,y=longput]{figdata/derivatives/01-no-arbitrage-and-the-fundamental-theorems/box_payoff.csv};
\addplot[omProp,thick,dashed] table[x=s,y=shortput]{figdata/derivatives/01-no-arbitrage-and-the-fundamental-theorems/box_payoff.csv};
\addplot[omThm,ultra thick] table[x=s,y=total]{figdata/derivatives/01-no-arbitrage-and-the-fundamental-theorems/box_payoff.csv};
\legend{call 5000 acheté,call 6000 vendu,put 6000 acheté,put 5000 vendu,le box}
\end{axis}
\end{tikzpicture}

\omcaption{Les quatre jambes d'un \omterm{def:dv:no-arbitrage-and-the-fundamental-theorems:box}{box spread} 5\,000--6\,000 à l'échéance. Chaque jambe présente
son point de rupture à son strike~; les points de rupture des calls et ceux des puts se compensent,
et la somme est la droite horizontale à 1\,000 points. Données~: le code
du chapitre.}\label{fig:dv:no-arbitrage-and-the-fundamental-theorems:box}
\end{omfigure}

\begin{example}[Lire le taux sur un box]\label{ex:dv:no-arbitrage-and-the-fundamental-theorems:boxrate}
Le box d'un an 5\,000--6\,000 de l'introduction se traite à 958.90 points. Son
taux à capitalisation continue vaut $-\ln(958.90/1000)=4.20\%$. Un prêteur qui l'achète
débourse USD~95\,890 et reçoit USD~100\,000 dans un an (le multiplicateur est
USD~100 par point)~; un emprunteur le vend.
\end{example}

\begin{dated}{2026-09}{Le marché des box spreads sur indice}\label{dat:dv:no-arbitrage-and-the-fundamental-theorems:box}
La bourse qui cote les options sur l'indice S\&P~500 publie un volume quotidien moyen de
USD~929 millions de notionnel en \omterm{def:dv:no-arbitrage-and-the-fundamental-theorems:box}{box spreads} sur indice en 2024. Les box sont en majorité
construits avec des écarts de strikes de 1\,000 points, un notionnel de USD~100\,000 par box, et la
bourse indique que leurs taux se situent généralement entre 10 et 55 points de base au-dessus
des rendements des bons du Trésor à trois mois. Chaque jambe est compensée, si bien que la
contrepartie du prêteur est la chambre de compensation des options, et non l'emprunteur.
\end{dated}

\section{Prix d'état sur une période}

\begin{definition}[Titre d'Arrow--Debreu et prix d'état]\label{def:dv:no-arbitrage-and-the-fundamental-theorems:ad}
Un \emph{titre d'Arrow--Debreu}\index{titre d'Arrow--Debreu} (Arrow--Debreu security) de l'état $s$
verse 1 si l'état $s$ se réalise et rien sinon. Son prix $q_s$ est le
\emph{prix d'état}\index{prix d'état} (state price) de $s$. Un vecteur $q\in\R^S$ est un
vecteur de prix d'état du marché s'il évalue chaque actif négocié~:
$p = D^\top q$, soit $p_j=\sum_s q_sD_{sj}$.
\end{definition}

Si les \omterm{def:dv:no-arbitrage-and-the-fundamental-theorems:ad}{titres d'Arrow--Debreu} étaient négociés, toute créance $g$ serait le portefeuille
composé de $\sum_s g_s$ unités du titre $s$, et coûterait $q\cdot g$. Ils ne se négocient pas,
mais les \omterm{def:dv:no-arbitrage-and-the-fundamental-theorems:ad}{prix d'état} peuvent néanmoins exister~: ce sont les inconnues du système linéaire
$D^\top q=p$, une équation par actif.

\begin{theorem}[Absence d'arbitrage et prix d'état strictement positifs]\label{thm:dv:no-arbitrage-and-the-fundamental-theorems:ftap1}
Dans le modèle fini à une période, il n'existe pas d'\omterm{def:dv:no-arbitrage-and-the-fundamental-theorems:arbitrage}{arbitrage} si et seulement s'il existe un
vecteur de \omterm{def:dv:no-arbitrage-and-the-fundamental-theorems:ad}{prix d'état} avec $q_s>0$ pour chaque état.
\end{theorem}
\begin{proof}
Si $q>0$ évalue chaque actif et si $D\theta\ge0$, alors $p\cdot\theta=q\cdot
D\theta\ge0$, avec égalité seulement si $D\theta=0$~: aucun portefeuille n'est
un \omterm{def:dv:no-arbitrage-and-the-fundamental-theorems:arbitrage}{arbitrage}. Réciproquement, supposons qu'il n'existe pas d'\omterm{def:dv:no-arbitrage-and-the-fundamental-theorems:arbitrage}{arbitrage}. L'ensemble
$C=\{(-p\cdot\theta, D\theta):\theta\in\R^N\}$ est un sous-espace vectoriel de
$\R^{1+S}$ qui ne rencontre l'orthant positif qu'en zéro. Par le théorème de
séparation des hyperplans (lemme de Farkas, livre~4, chapitre~23), il existe un
vecteur $(\phi_0,\phi)$ à coordonnées strictement positives orthogonal à $C$~:
$-\phi_0\,p\cdot\theta+\phi\cdot D\theta=0$ pour tout $\theta$. Alors
$q=\phi/\phi_0>0$ vérifie $p=D^\top q$.
\end{proof}

Les \omterm{def:dv:no-arbitrage-and-the-fundamental-theorems:ad}{prix d'état} recouvrent deux objets familiers. Si une obligation sans risque versant 1 dans
chaque état est négociée, son prix vaut $\sum_sq_s=P(0,T)$, et les
$\pi_s=q_s/P(0,T)$ sont, une fois normalisés, des poids positifs de somme égale à un~: une probabilité, sous laquelle
le prix de chaque actif est son espérance de gain actualisée,
$p_j=P(0,T)\sum_s\pi_sD_{sj}$. Cette probabilité est la probabilité risque-neutre du
livre~4, chapitre~5, la probabilité martingale équivalente, le compte bancaire servant de
numéraire~: les prix divisés par l'obligation sont des martingales. Rien dans cette
construction n'indique que les états soient probables ou improbables~; c'est une normalisation des prix.

\begin{definition}[Densité des prix d'état]\label{def:dv:no-arbitrage-and-the-fundamental-theorems:spd}
Étant données les probabilités réelles $\P(s)>0$ des états, la
\emph{densité des prix d'état}\index{densité des prix d'état} (state-price density) est $m_s=q_s/\P(s)$~; chaque
prix est une espérance sous $\P$~: $p_j=\E[mD_j]$. C'est le rapport entre la
probabilité risque-neutre et la probabilité réelle, multiplié par le facteur d'actualisation, et elle est
élevée dans les états où les investisseurs valorisent le plus un gain.
\end{definition}

\begin{example}[Un marché binomial]\label{ex:dv:no-arbitrage-and-the-fundamental-theorems:binomial}
Une obligation coûte 0.98 et verse 1~; une action coûte 100 et verse 120 ou 90. Les
équations $q_u+q_d=0.98$ et $120q_u+90q_d=100$ donnent $q_u=0.3933$, $q_d=0.5867$,
tous deux positifs~: pas d'\omterm{def:dv:no-arbitrage-and-the-fundamental-theorems:arbitrage}{arbitrage}. La probabilité risque-neutre de l'état haussier vaut
$0.3933/0.98=0.4014$, quelle que soit sa probabilité réelle. Un call de strike 100
verse $(20,0)$ et vaut $20q_u=7.8667$.
\end{example}

\section{Mesures martingales et premier théorème fondamental}

Sur plusieurs périodes, un portefeuille est rééquilibré à chaque date $t_k$ en fonction
de ce qui est connu à cet instant. Soit $S_t$ le vecteur des prix risqués, $B_t$ le compte
bancaire, $\theta_k$ les positions choisies en $t_k$ et conservées jusqu'en $t_{k+1}$.

\begin{definition}[Stratégie autofinancée]\label{def:dv:no-arbitrage-and-the-fundamental-theorems:sf}
Une stratégie de trading $(\theta_k)$, adaptée à l'information disponible à chaque
date, est une \emph{stratégie autofinancée}\index{stratégie autofinancée} (self-financing strategy) si
rien n'est ajouté ni retiré entre $0$ et $T$~: chaque rééquilibrage est financé
par le portefeuille lui-même, de sorte que sa valeur ne varie que par les variations de prix,
$V_{t_{k+1}}-V_{t_k}=\theta_k\cdot(S_{t_{k+1}}-S_{t_k})$ en unités actualisées. En
temps continu, $dV_t=\theta_t\,dS_t$ à l'aide d'une intégrale stochastique.
\end{definition}

Un \omterm{def:dv:no-arbitrage-and-the-fundamental-theorems:arbitrage}{arbitrage} sur plusieurs périodes est une \omterm{def:dv:no-arbitrage-and-the-fundamental-theorems:sf}{stratégie autofinancée} avec $V_0\le0$,
$V_T\ge0$ et $\P(V_T>0)>0$ (ou $V_0<0$). Le résultat à une période s'étend date
par date.

\begin{theorem}[Premier théorème fondamental de l'évaluation des actifs]\label{thm:dv:no-arbitrage-and-the-fundamental-theorems:ftap}
Dans un marché comportant un nombre fini de dates de négociation, le \emph{théorème fondamental de
l'évaluation des actifs}\index{théorème fondamental de l'évaluation des actifs} (fundamental theorem of asset pricing) affirme qu'il n'existe pas
d'\omterm{def:dv:no-arbitrage-and-the-fundamental-theorems:arbitrage}{arbitrage} si et seulement s'il existe une probabilité martingale équivalente~: une
probabilité $\mathbb Q$ équivalente à $\P$ sous laquelle chaque prix actualisé $S_t/B_t$ est une
martingale. Le prix sans \omterm{def:dv:no-arbitrage-and-the-fundamental-theorems:arbitrage}{arbitrage} d'une créance $g$ versée en $T$
vaut alors $B_0\E^{\mathbb Q}[g/B_T]$ pour une telle $\mathbb Q$.
\end{theorem}
\admitted

Pour un espace d'états fini, la preuve est le théorème à une période appliqué à chaque nœud de
l'arbre d'information, les \omterm{def:dv:no-arbitrage-and-the-fundamental-theorems:ad}{prix d'état} de chaque nœud donnant les probabilités conditionnelles de
$\mathbb Q$ (Harrison et Kreps, puis Harrison et
Pliska). Pour le temps discret général, il est dû à Dalang, Morton et Willinger~;
en temps continu, «~absence d'\omterm{def:dv:no-arbitrage-and-the-fundamental-theorems:arbitrage}{arbitrage}~» doit être renforcée en «~absence de repas gratuit
à risque évanescent~» et la propriété de martingale affaiblie en sigma-martingale
(Delbaen et Schachermayer)~; les énoncés de théorie de la mesure et le théorème de Girsanov
qui construit $\mathbb Q$ dans les modèles de diffusion, figurent dans le livre~4,
aux chapitres~1 et~5. Ce dont le desk a besoin, c'est le sens d'utilisation~: pour évaluer, il
faut trouver une mesure sous laquelle les prix négociés actualisés sont des martingales, puis
en calculer l'espérance.

\begin{remark}[Pourquoi le compte bancaire]\label{rem:dv:no-arbitrage-and-the-fundamental-theorems:numeraire}
Rien n'impose le compte bancaire comme unité. Tout actif négocié à prix positif peut servir de
numéraire $\mathcal N_t$~: sous sa mesure $\mathbb Q^{\mathcal N}$, les prix divisés par $\mathcal N_t$
sont des martingales. Bien choisir le numéraire
(l'obligation zéro-coupon pour un gain en $T$, l'actif lui-même pour un gain en
actions) supprime un facteur d'actualisation stochastique d'une espérance~; le chapitre~3 l'utilise pour
dériver la formule de Black--Scholes en deux lignes.
\end{remark}

\section{Complétude, réplication et second théorème}

\begin{definition}[Marché complet]\label{def:dv:no-arbitrage-and-the-fundamental-theorems:complete}
Un \emph{marché complet}\index{marché complet} (complete market) est un marché dans lequel toute
créance contingente admet une \omterm{def:dv:no-arbitrage-and-the-fundamental-theorems:sf}{stratégie autofinancée} de réplication. Dans le modèle fini à une
période, le marché est complet si et seulement si $D$ est de rang $S$.
\end{definition}

\begin{theorem}[Second théorème fondamental]\label{thm:dv:no-arbitrage-and-the-fundamental-theorems:ftap2}
Dans un marché sans \omterm{def:dv:no-arbitrage-and-the-fundamental-theorems:arbitrage}{arbitrage} à dates et états en nombre fini, le marché est
complet si et seulement si la probabilité martingale équivalente est unique.
\end{theorem}
\begin{proof}[Preuve partielle (une période)]
Les vecteurs de \omterm{def:dv:no-arbitrage-and-the-fundamental-theorems:ad}{prix d'état} sont les solutions de $D^\top q=p$ avec $q>0$. Si $D$
est de rang $S$, $D^\top$ est injective et la solution est unique. Si le rang est
strictement inférieur à $S$, le noyau de $D^\top$ contient un $z\ne0$, et $q+\varepsilon z$
reste une solution positive pour $\varepsilon$ petit~: la mesure n'est pas unique.
\end{proof}

Un \omterm{def:dv:no-arbitrage-and-the-fundamental-theorems:complete}{marché complet} valorise toute créance~: le prix est le coût du \omterm{def:dv:no-arbitrage-and-the-fundamental-theorems:claim}{portefeuille de réplication},
et détenir ce portefeuille face à la vente de la créance ne laisse aucun risque. C'est ce qui
fait de la valorisation d'une option une transaction et non une opinion. Le
marché binomial de l'\cref{ex:dv:no-arbitrage-and-the-fundamental-theorems:binomial}
est complet~; le chapitre~2 en construit toute une théorie de la valorisation. Si l'on ajoute
un troisième état, il ne l'est plus.

\begin{omfigure}
\begin{tikzpicture}[font=\small,
  nd/.style={draw=omDef,thick,rounded corners=2pt,fill=omDef!6,align=center,minimum width=2.4cm,inner sep=3pt}]
\node[nd] (o) at (0,0) {aujourd'hui\\obligation 0.98, action 100};
\node[nd] (u) at (6,1.7) {état $u$~: action 120};
\node[nd] (m) at (6,0) {état $m$~: action 100};
\node[nd] (d) at (6,-1.7) {état $d$~: action 80};
\draw[-{Stealth},thick,gray] (o.east) -- (u.west);
\draw[-{Stealth},thick,gray] (o.east) -- (m.west);
\draw[-{Stealth},thick,gray] (o.east) -- (d.west);
\node[anchor=west,align=left,font=\footnotesize] at (8.3,1.7) {$q_u\in(0.10,\,0.54)$};
\node[anchor=west,align=left,font=\footnotesize] at (8.3,0) {$q_m\in(0,\,0.88)$};
\node[anchor=west,align=left,font=\footnotesize] at (8.3,-1.7) {$q_d\in(0,\,0.44)$};
\node[align=center,font=\footnotesize,text=omThm] at (3,-2.6) {sommets $(0.10,\,0.88,\,0)$ et $(0.54,\,0,\,0.44)$~:\\tout vecteur de prix d'état positif se situe strictement entre eux};
\end{tikzpicture}

\omcaption{Le marché trinomial~: deux actifs, trois états. Les deux équations
$q_u+q_m+q_d=0.98$ et $120q_u+100q_m+80q_d=100$ laissent un degré de liberté, un
segment de vecteurs de \omterm{def:dv:no-arbitrage-and-the-fundamental-theorems:ad}{prix d'état} dont chaque extrémité annule le prix d'un état.
Nombres illustratifs.}\label{fig:dv:no-arbitrage-and-the-fundamental-theorems:tree}
\end{omfigure}

\section{Bornes de prix lorsque le marché est incomplet}

Dans le marché trinomial de la \cref{fig:dv:no-arbitrage-and-the-fundamental-theorems:tree},
un call de strike 100 paie $(20,0,0)$. Son prix doit être $20q_u$ pour un certain
vecteur de \omterm{def:dv:no-arbitrage-and-the-fundamental-theorems:ad}{prix d'état} positif, donc tout montant strictement compris entre $20\times0.10=2$ et
$20\times0.54=10.80$ est compatible avec l'absence d'\omterm{def:dv:no-arbitrage-and-the-fundamental-theorems:arbitrage}{arbitrage}. Le marché n'a pas tranché.

\begin{definition}[Prix de sur-réplication]\label{def:dv:no-arbitrage-and-the-fundamental-theorems:super}
Le \emph{prix de sur-réplication}\index{prix de sur-réplication} (super-replication price) d'une créance
$g$ est le coût minimal d'un portefeuille d'actifs négociés qui paie au moins $g$ dans
chaque état~: $\inf\{p\cdot\theta: D\theta\ge g\}$. Le prix de sous-réplication est le
coût maximal d'un portefeuille qui paie au plus $g$.
\end{definition}

\begin{proposition}[L'intervalle sans arbitrage]\label{prop:dv:no-arbitrage-and-the-fundamental-theorems:interval}
Dans un marché sans \omterm{def:dv:no-arbitrage-and-the-fundamental-theorems:arbitrage}{arbitrage} à une période, le \omterm{def:dv:no-arbitrage-and-the-fundamental-theorems:super}{prix de sur-réplication} de $g$ est égal à
$\max q\cdot g$ sur les vecteurs de \omterm{def:dv:no-arbitrage-and-the-fundamental-theorems:ad}{prix d'état} $q\ge0$ tels que $D^\top q=p$, et le
prix de sous-réplication est égal au minimum. Un nouvel actif payant $g$ peut être ajouté au
prix $c$ sans créer d'\omterm{def:dv:no-arbitrage-and-the-fundamental-theorems:arbitrage}{arbitrage} si et seulement si $c$ se situe strictement entre
eux, ou bien $g$ est réplicable et $c$ est son prix.
\end{proposition}
\begin{proof}
Le premier énoncé relève de la dualité des programmes linéaires~: le dual de $\min\{p\cdot
\theta:D\theta\ge g\}$ est $\max\{q\cdot g:D^\top q=p,\ q\ge0\}$, et les deux optima
sont atteints et égaux lorsque les ensembles admissibles sont non vides. Pour le second,
appliquer le \cref{thm:dv:no-arbitrage-and-the-fundamental-theorems:ftap1} au
marché élargi~: un vecteur de \omterm{def:dv:no-arbitrage-and-the-fundamental-theorems:ad}{prix d'état} positif existe pour celui-ci exactement lorsque
$c=q\cdot g$ pour un certain $q>0$ évaluant les anciens actifs, et les valeurs $q\cdot g$ prises
sur l'ouvert des $q$ positifs remplissent l'intervalle ouvert (ou un point unique lorsque $g$
est réplicable).
\end{proof}

À l'extrémité de l'intervalle, la créance constitue encore un \omterm{def:dv:no-arbitrage-and-the-fundamental-theorems:arbitrage}{arbitrage}~: à 10.80, le call
peut être vendu et sur-répliqué à coût net nul par $0.5$ action et $-40$
obligations, et la position couverte paie $(0,10,0)$, rien dans deux états et 10 dans
le troisième. La \Cref{fig:dv:no-arbitrage-and-the-fundamental-theorems:bounds} montre
l'intervalle par strike~: large à la monnaie, là où le paiement du call dépend de
l'état que les deux actifs ne peuvent distinguer, refermé aux extrémités de la
plage de strikes, là où le call est réplicable.

\begin{omfigure}
\begin{tikzpicture}
\begin{axis}[width=10.5cm,height=5cm,
  xlabel={strike},ylabel={prix du call},
  tick label style={font=\small},label style={font=\small},
  xmin=80,xmax=120,ymin=0,ymax=23,grid=major,grid style={gray!25},
  legend columns=3,legend style={font=\footnotesize,at={(0.5,-0.3)},anchor=north,draw=none,fill=none,/tikz/every even column/.append style={column sep=8pt}},
  table/col sep=comma]
\addplot[name path=up,omThm,thick] table[x=strike,y=upper]{figdata/derivatives/01-no-arbitrage-and-the-fundamental-theorems/bounds.csv};
\addplot[name path=lo,omDef,thick] table[x=strike,y=lower]{figdata/derivatives/01-no-arbitrage-and-the-fundamental-theorems/bounds.csv};
\addplot[omProp!25] fill between[of=up and lo];
\addplot[omProp,thick,dashed] table[x=strike,y=complete]{figdata/derivatives/01-no-arbitrage-and-the-fundamental-theorems/bounds.csv};
\legend{sur-réplication,sous-réplication,,une fois le call 100 échangé à 6}
\end{axis}
\end{tikzpicture}

\omcaption{L'intervalle sans \omterm{def:dv:no-arbitrage-and-the-fundamental-theorems:arbitrage}{arbitrage} d'un call dans le marché trinomial, par
strike. Dès qu'une option est échangée (le call 100 à 6.00), le marché est complet et
chaque autre strike possède un prix unique, la ligne en tirets, à l'intérieur de la bande. Données~:
le tutoriel.}\label{fig:dv:no-arbitrage-and-the-fundamental-theorems:bounds}
\end{omfigure}

\begin{method}[Bornes statiques sur une chaîne cotée]\label{met:dv:no-arbitrage-and-the-fundamental-theorems:chain}
Pour des calls européens sur un même sous-jacent de forward $F$ et de facteur d'actualisation $P$,
sans aucun modèle~:
\begin{enumerate}
\item $P\max(F-K,0)\le C(K)\le PF$~; pour les puts, $P\max(K-F,0)\le P(K)\le PK$~;
\item $C$ décroît en $K$ et $C(K_1)-C(K_2)\le P(K_2-K_1)$ (le spread de calls
paie au plus l'écart de strikes)~;
\item $C$ est convexe en $K$~: un butterfly coûte au moins zéro~;
\item $C(K)-P(K)=P(F-K)$ à chaque strike, et chaque box coûte $P(K_2-K_1)$.
\end{enumerate}
Lorsque les cotations comportent un cours acheteur et un cours vendeur, chaque contrôle est effectué sur des prix exécutables~; un échec est un
\omterm{def:dv:no-arbitrage-and-the-fundamental-theorems:arbitrage}{arbitrage} dont le portefeuille est désigné par le contrôle
(\cref{fig:dv:no-arbitrage-and-the-fundamental-theorems:region}).
\end{method}

\begin{omfigure}
\begin{tikzpicture}
\begin{axis}[width=10.5cm,height=5cm,
  xlabel={strike},ylabel={prix du call (points d'indice)},
  tick label style={font=\small},label style={font=\small},
  xmin=3000,xmax=8000,ymin=0,ymax=6200,grid=major,grid style={gray!25},xtick={3000,4000,5000,6000,7000,8000},
  xticklabel style={/pgf/number format/.cd,1000 sep={\,}},yticklabel style={/pgf/number format/.cd,1000 sep={\,}},
  legend columns=3,legend style={font=\footnotesize,at={(0.5,-0.3)},anchor=north,draw=none,fill=none,/tikz/every even column/.append style={column sep=8pt}},
  table/col sep=comma]
\addplot[omThm,thick,dashed] table[x=strike,y=upper]{figdata/derivatives/01-no-arbitrage-and-the-fundamental-theorems/call_region.csv};
\addplot[omDef,thick,dashed] table[x=strike,y=lower]{figdata/derivatives/01-no-arbitrage-and-the-fundamental-theorems/call_region.csv};
\addplot[omProp,very thick] table[x=strike,y=call]{figdata/derivatives/01-no-arbitrage-and-the-fundamental-theorems/call_region.csv};
\legend{forward actualisé,valeur intrinsèque actualisée,courbe de call cohérente avec le smile}
\end{axis}
\end{tikzpicture}

\omcaption{Où peut se situer un prix de call~: entre la valeur intrinsèque actualisée
$P\max(F-K,0)$ et le forward actualisé $PF$, décroissant et convexe en fonction
du strike. La courbe pleine est la chaîne du problème du week-end (forward 5\,850,
un an, smile dissymétrique). Paramètres
illustratifs.}\label{fig:dv:no-arbitrage-and-the-fundamental-theorems:region}
\end{omfigure}

\section{Tutoriel~: prix d'état, bornes et taux d'un box-spread}\label{tut:dv:no-arbitrage-and-the-fundamental-theorems:tut}

\begin{tutorial}
\textbf{Objectif.} Résoudre les \omterm{def:dv:no-arbitrage-and-the-fundamental-theorems:ad}{prix d'état}, calculer l'intervalle d'absence d'\omterm{def:dv:no-arbitrage-and-the-fundamental-theorems:arbitrage}{arbitrage}
d'un \omterm{def:dv:no-arbitrage-and-the-fundamental-theorems:claim}{actif contingent} sur un marché incomplet, trouver un portefeuille d'\omterm{def:dv:no-arbitrage-and-the-fundamental-theorems:arbitrage}{arbitrage} lorsqu'un prix
sort de cet intervalle, et lire un taux sur un box coté. \textbf{État final~:}
la \cref{fig:dv:no-arbitrage-and-the-fundamental-theorems:bounds} et les valeurs de
l'\cref{ex:dv:no-arbitrage-and-the-fundamental-theorems:binomial} et celles du
problème du week-end.
\begin{tutsteps}
\item \textbf{Sommets de l'ensemble des \omterm{def:dv:no-arbitrage-and-the-fundamental-theorems:ad}{prix d'état}.} Chaque sommet résout le système sur
un support comportant autant d'états que d'actifs~; les bornes d'un \omterm{def:dv:no-arbitrage-and-the-fundamental-theorems:claim}{actif contingent} sont
atteintes aux sommets.
\omcode{code/derivatives/01-no-arbitrage-and-the-fundamental-theorems/python/dv_arbitrage.py}{33}{52}{Sommets des prix d'état et bornes d'absence d'arbitrage d'un actif contingent.}\label{lst:dv:no-arbitrage-and-the-fundamental-theorems:vertices}
\item \textbf{Le côté primal.} Le portefeuille de sur-réplication le moins cher, et une
recherche de portefeuille d'\omterm{def:dv:no-arbitrage-and-the-fundamental-theorems:arbitrage}{arbitrage}, sont des programmes linéaires~; la fonction
\texttt{lp\_max} de la construction les résout exactement en énumérant les solutions de base.
\omcode{code/derivatives/01-no-arbitrage-and-the-fundamental-theorems/python/dv_arbitrage.py}{55}{74}{Sur-réplication et recherche d'arbitrage sous forme de programmes linéaires.}\label{lst:dv:no-arbitrage-and-the-fundamental-theorems:lp}
\item \textbf{Exécuter} les tests~: la borne supérieure égale le coût de sur-réplication
pour chaque strike, le call spread 90--110 a un intervalle d'un point, et pricer
le call 100 à 11.00 renvoie l'\omterm{def:dv:no-arbitrage-and-the-fundamental-theorems:arbitrage}{arbitrage} $(-40, 0.5, -1)$ avec un edge de 0.20.
\item \textbf{Le box.} Pricer les quatre jambes à des prix exécutables puis inverser~:
\omcode{code/firm/arbcheck/firm_arbcheck.py}{92}{102}{Le coût d'achat ou de vente d'un box, et le taux qu'il implique.}\label{lst:dv:no-arbitrage-and-the-fundamental-theorems:box}
\end{tutsteps}
\textbf{Pour aller plus loin.} Ajouter un quatrième état et une seconde option, et observer
l'intervalle d'une troisième option se resserrer~; puis coter les quatre jambes du box
d'un tick plus large et voir de combien les taux de prêt et d'emprunt s'écartent.
\end{tutorial}

\section{Construction~: le contrôleur d'arbitrage statique}\label{bld:dv:no-arbitrage-and-the-fundamental-theorems:arbcheck}

\begin{build}
\bfield{Rôle} Avant qu'un modèle ne touche une chaîne de cotations, la firme miniature
vérifie que la chaîne elle-même est exempte d'\omterm{def:dv:no-arbitrage-and-the-fundamental-theorems:arbitrage}{arbitrage} statique, et signale le
portefeuille et l'edge de chaque violation. Les mêmes contrôles protègent les ajustements
de surface des chapitres 7 et 8.
\bfield{Interface} \texttt{Quote(strike, right, bid, ask)}~;
\texttt{check\_chain(quotes, forward, df) -> list[Violation]} avec
\texttt{Violation(kind, strikes, edge, legs)}~; \texttt{box\_cost},
\texttt{box\_rate}, \texttt{check\_box}~; \texttt{lp\_max(c, a\_ub, b\_ub)}.
\bfield{Règles} Acheter au cours vendeur, vendre au cours acheteur~; un edge est un encaissement immédiat
d'un portefeuille dont le paiement n'est jamais négatif~; la convexité se contrôle sur des
écarts de strikes inégaux, avec les poids $(K_3-K_2)/(K_3-K_1)$ et $(K_2-K_1)/(K_3-K_1)$.
\bfield{Tests d'acceptation} \texttt{code/firm/arbcheck/tests/}~: une chaîne de
Black--Scholes passe~; un strike mal pricé est repéré par le contrôle de convexité avec un
edge positif~; une option inférieure à sa valeur intrinsèque actualisée est repérée avec le
bon edge~; le taux du box d'une chaîne propre égale le taux de la chaîne.
\bfield{Pour aller plus loin} Remplacer la liste de règles par un seul programme linéaire
sur tous les portefeuilles statiques de la chaîne (options, forward, obligation), qui trouve
aussi les violations combinant plus de trois strikes.
\end{build}

\begin{omsources}
\item J.~M.~Harrison et D.~M.~Kreps, «~Martingales and arbitrage in multiperiod
securities markets~», \emph{Journal of Economic Theory} 20 (1979) 381--408.
\item J.~M.~Harrison et S.~R.~Pliska, «~Martingales and stochastic integrals in
the theory of continuous trading~», \emph{Stochastic Processes and their
Applications} 11 (1981) 215--260.
\item F.~Delbaen et W.~Schachermayer, «~A general version of the fundamental
theorem of asset pricing~», \emph{Mathematische Annalen} 300 (1994) 463--520.
\item K.~J.~Arrow, «~The role of securities in the optimal allocation of
risk-bearing~», \emph{Review of Economic Studies} 31 (1964) 91--96.
\item J.~H.~van Binsbergen, W.~F.~Diamond et M.~Grotteria, «~Risk-free
interest rates~», \emph{Journal of Financial Economics} 143 (2022) 1--29.
\item Cboe, SPX box spreads (documentation produit, 2024--2025).
\end{omsources}

\section{Exercices}

\begin{exercise}[$\star$]\label{exo:dv:no-arbitrage-and-the-fundamental-theorems:1}
Un zéro-coupon coûte 0.99 et paie 1~; une action coûte 100 et paie 115 ou 95. Calculer les
\omterm{def:dv:no-arbitrage-and-the-fundamental-theorems:ad}{prix d'état}, la probabilité risque-neutre de l'état haussier, et le prix d'un call de
strike 100.
\end{exercise}

\begin{exercise}[$\star$]\label{exo:dv:no-arbitrage-and-the-fundamental-theorems:2}
Vérifier, pour un indice à l'échéance inférieure à 5\,000, entre les strikes, et supérieure
à 6\,000, que le box 5\,000--6\,000 paie 1\,000 points d'indice.
\end{exercise}

\begin{exercise}[$\star$]\label{exo:dv:no-arbitrage-and-the-fundamental-theorems:3}
Dans le marché trinomial, un put de strike 100 est offert à 9.00. Y a-t-il un
\omterm{def:dv:no-arbitrage-and-the-fundamental-theorems:arbitrage}{arbitrage}~? Si oui, donner le portefeuille et son edge.
\end{exercise}

\begin{exercise}[$\star\star$]\label{exo:dv:no-arbitrage-and-the-fundamental-theorems:4}
Construire un marché à deux états satisfaisant la \omterm{def:dv:no-arbitrage-and-the-fundamental-theorems:loop}{loi du prix unique} et admettant un
\omterm{def:dv:no-arbitrage-and-the-fundamental-theorems:arbitrage}{arbitrage}. Laquelle des conditions du
\cref{thm:dv:no-arbitrage-and-the-fundamental-theorems:ftap1} n'est pas vérifiée~?
\end{exercise}

\begin{exercise}[$\star\star$]\label{exo:dv:no-arbitrage-and-the-fundamental-theorems:5}
Dans le marché trinomial complété par le call de strike 100 à 6.00, les \omterm{def:dv:no-arbitrage-and-the-fundamental-theorems:ad}{prix d'état} sont
$(0.30, 0.48, 0.20)$. Avec les probabilités du monde réel $(0.45, 0.45, 0.10)$, donner
les probabilités risque-neutres et la \omterm{def:dv:no-arbitrage-and-the-fundamental-theorems:spd}{densité des prix d'état}, et dire ce que la
forme de la densité signifie.
\end{exercise}

\begin{exercise}[$\star\star$]\label{exo:dv:no-arbitrage-and-the-fundamental-theorems:6}
Donner l'intervalle d'absence d'\omterm{def:dv:no-arbitrage-and-the-fundamental-theorems:arbitrage}{arbitrage} d'un digital qui verse 1 dans l'état haussier du
marché trinomial, et le portefeuille de sur-réplication à son extrémité supérieure.
\end{exercise}

\begin{exercise}[$\star\star\star$]\label{exo:dv:no-arbitrage-and-the-fundamental-theorems:7}
\emph{Programmation.} Avec \texttt{price\_bounds}, calculer l'intervalle de
le call spread 90--110 dans le marché trinomial. Expliquer le résultat.
\end{exercise}

\begin{exercise}[$\star\star\star$]\label{exo:dv:no-arbitrage-and-the-fundamental-theorems:8}
\emph{Repérer l'erreur.} «~Dans le marché trinomial, le call 100 peut s'échanger
entre 2 et 10.80. Mon modèle l'évalue à 6, donc je l'achète à 5 et je me couvre~: cela
constitue un profit sans risque de 1.~»~Corriger le raisonnement.
\end{exercise}

\section{Problème~: Taux du box spread}

\begin{problem}[{Problème du week-end --- prêter des liquidités au moyen de quatre options}]\label{pb:dv:no-arbitrage-and-the-fundamental-theorems:1}
Une trésorière souhaite prêter 10~millions de dollars pendant un an au moyen de
\omterm{def:dv:no-arbitrage-and-the-fundamental-theorems:box}{box spreads} sur indices et compare les cotations au taux overnight, qu'elle
prévoit de voir capitaliser à 4.00~\% (en capitalisation continue) sur l'année. Les options à un an sur
l'indice, européennes et réglées en espèces avec un multiplicateur de 100~dollars, sont cotées (en
points d'indice)~: 5\,000 call 971.80--973.40, 6\,000 call 324.60--326.20, 5\,000 put
156.70--158.30, 6\,000 put 468.40--470.00. Les cotations sont données à titre illustratif.

\medskip\noindent\textbf{Partie I --- Les quatre jambes.}
\begin{enumerate}
\item Écrire le gain de chaque jambe à l'échéance et montrer que le total vaut 1\,000
points dans chaque état.
\item Combien verse un box en dollars~?
\item Combien la trésorière paie-t-elle pour un box acheté jambe par jambe aux cotations~?
\item Que recevrait un emprunteur pour un box vendu jambe par jambe~?
\item Quels risques le box élimine-t-il, et lesquels subsistent-ils~?
\end{enumerate}

\medskip\noindent\textbf{Partie II --- Taux.}
\begin{enumerate}[resume]
\item Donner le taux de prêt implicite dans l'achat aux cotations.
\item Donner le taux d'emprunt implicite dans la vente aux cotations.
\item Donner le prix milieu du box et son taux.
\item Combien de dollars d'intérêts la trésorière gagne-t-elle par box acheté jambe par jambe~?
\item Donner les spreads des taux de prêt et d'emprunt par rapport au taux overnight, en points de base.
\end{enumerate}

\medskip\noindent\textbf{Partie III --- Paquets et \omterm{def:dv:no-arbitrage-and-the-fundamental-theorems:arbitrage}{arbitrage}.}
\begin{enumerate}[resume]
\item Pourquoi les box se négocient-ils en paquet à un seul prix plutôt que jambe par jambe~?
\item En prenant le taux overnight comme taux d'actualisation, existe-t-il un \omterm{def:dv:no-arbitrage-and-the-fundamental-theorems:arbitrage}{arbitrage} dans les cotations~?
\item À quel prix de paquet la trésorière gagnerait-elle exactement le taux overnight~?
\item Qu'est-ce qui changerait si les options étaient américaines, comme le sont les options sur actions~?
\item Donner une raison pour laquelle les taux des box peuvent dépasser les rendements des bons du Trésor (T-bills) de même maturité.
\end{enumerate}

\medskip\noindent\textbf{Partie IV --- Jugement.}
\begin{enumerate}[resume]
\item Qui emprunte en vendant des box, et pourquoi~?
\item Combien de box lui faut-il, et combien paie-t-elle au prix milieu~?
\item Que détient-elle si le membre compensateur qu'elle utilise fait défaut~?
\item Énoncer le \emph{résultat nommé}~: le taux implicite au prix milieu du box et son spread par rapport au taux overnight.
\item En une phrase~: pourquoi un \omterm{def:dv:no-arbitrage-and-the-fundamental-theorems:box}{box spread} est-il une obligation zéro-coupon~?
\end{enumerate}
\end{problem}

\section{Questions d'entretien}

\begin{interviewq}[$\star$ \iqroles{trader, chercheur}]\label{iq:dv:no-arbitrage-and-the-fundamental-theorems:1}
Qu'est-ce qu'un \omterm{def:dv:no-arbitrage-and-the-fundamental-theorems:box}{box spread}, quel est son paiement à l'échéance, et que vous apprend son prix~?
\end{interviewq}

\begin{interviewq}[$\star$ \iqroles{chercheur, banque}]\label{iq:dv:no-arbitrage-and-the-fundamental-theorems:2}
Énoncer le \omterm{thm:dv:no-arbitrage-and-the-fundamental-theorems:ftap}{théorème fondamental de l'évaluation des actifs}, et dire ce que sont et ce que ne sont pas les
probabilités risque-neutres.
\end{interviewq}

\begin{interviewq}[$\star\star$ \iqroles{trader}]\label{iq:dv:no-arbitrage-and-the-fundamental-theorems:3}
Une action vaut 100 et vaudra 110 ou 95 dans un an~; le taux à un an est
nul. Évaluer le prix d'un digital qui verse 1 si l'action finit à 110.
\end{interviewq}

\begin{interviewq}[$\star\star$ \iqroles{chercheur}]\label{iq:dv:no-arbitrage-and-the-fundamental-theorems:4}
Pourquoi un \omterm{def:dv:no-arbitrage-and-the-fundamental-theorems:complete}{marché complet} est-il équivalent à une unique probabilité martingale~? Donner un
exemple de marché incomplet et dire ce qui advient des prix des options dans celui-ci.
\end{interviewq}

\begin{interviewq}[$\star\star$ \iqroles{trader, risque}]\label{iq:dv:no-arbitrage-and-the-fundamental-theorems:5}
Les calls de strikes 90, 100 et 110 ont pour prix milieu 12.00, 6.50 et 0.90,
chacun coté avec une fourchette de 0.10. Existe-t-il un \omterm{def:dv:no-arbitrage-and-the-fundamental-theorems:arbitrage}{arbitrage}~? Que négocieriez-vous~?
\end{interviewq}

\begin{interviewq}[$\star\star\star$ \iqroles{développeur}]\label{iq:dv:no-arbitrage-and-the-fundamental-theorems:6}
Concevoir un service qui vérifie 3\,000 cotations d'options cotées sur un même
sous-jacent afin de détecter l'\omterm{def:dv:no-arbitrage-and-the-fundamental-theorems:arbitrage}{arbitrage} statique à chaque changement de cotation.
\end{interviewq}
